% ============================================================================
% SECTION 4 --- Methods
%   2.1 exact transition support (state space, A(x), executor, closure,
%       trans-dimensionality; ONE proposition; families -> app:operators)
%   2.2 learning which legal edits are plausible (generator matching; arrive
%       quickly at R_theta over canonical successors; source independence)
%   2.3 canonical molecular successors (worked example; why aggregation)
% Terminology contract: "connected, valence-admissible within the declared
% molecular state space", "reference edit law", "canonical molecular
% successor".  "exact" is used only of the executor's support and of canonical
% normalization.  No number appears in this section.
% COMPRESSED TO THE 9-PAGE LIMIT.  Detail relocated to app:operators and
% app:quotient, which are unlimited; do not re-expand here without cutting
% elsewhere.
% ============================================================================
\section{Methods}
\label{sec:process}

We develop \method{} as a reusable stochastic process over molecular edits and a finite-horizon controller over that process. \method{} realizes generative stochastic rewriting by repeatedly applying executable molecular edits under a learned, goal-independent transition law. An exact rewrite system first determines which transitions between complete molecular graphs are executable. We then learn a goal-independent reference law over this state-dependent support and push it forward to a canonical molecular successor kernel $R_\theta$. Finally, we control the frozen successor process through a budget-conditioned $h$-transform, using a learned future-value model $h_\phi$ to approximate the required backward values at molecular scale.

\subsection{Executable molecular transitions}
\label{sec:support}

\method{} represents every generative state as a complete molecular graph. Let $n_{\max}$ denote the maximum number of heavy atoms admitted by the declared molecular state space. For $n\in\{1,\dots,n_{\max}\}$, let $\X_n$ be the set of connected molecular graphs with $n$ heavy atoms whose node and edge attributes lie in the declared vocabularies of elements, bond orders, formal charges, and implicit hydrogens, and in which every atom satisfies its declared valence constraints. The molecular state space is
\begin{equation}
\X=\bigsqcup_{n=1}^{n_{\max}}\X_n.
\label{eq:statespace}
\end{equation}
Thus each state $x\in\X$ is a complete molecule rather than a masked, corrupted, or partially specified representation.

Molecular transitions are generated by typed graph rewrites. An edit mark $a$ specifies a rewrite type together with its molecular match and typed payload, and $\T(x,a)$ denotes the graph obtained by applying that rewrite to $x$. The full atom-, bond-, and ring-level edit language is specified in \cref{app:operators}. From a state $x$, the executor exposes only the finite, state-dependent set
\begin{equation}
\A(x)=\bigl\{\,a:\ a\ \text{matches}\ x\ \ \text{and}\ \ \T(x,a)\ \text{passes all executor guards}\,\bigr\}.
\label{eq:editset}
\end{equation}
Execution follows a propose--validate--commit rule: a successor is admitted only if its attributes remain in the declared vocabularies, affected valences are satisfied, the graph remains connected, and its heavy-atom count lies in $\{1,\dots,n_{\max}\}$. An edit that fails any guard contributes no transition to $\A(x)$. Inadmissible molecular states are therefore absent from the stochastic support rather than discouraged through a learned penalty.

The decomposition of $\X$ by heavy-atom count makes the process trans-dimensional. Each admitted edit changes heavy-atom count by at most one, yielding \emph{birth}, \emph{death}, or \emph{same-cardinality} transitions between $\X_n$ and $\X_{n+1}$, $\X_{n-1}$, or $\X_n$, respectively. Birth and death edits create or remove typed graph entities that participate in valence, connectivity, and ring perception rather than merely activating or deactivating padded coordinates. A trajectory may therefore grow, shrink, and restructure a molecular graph within a common state space.

\begin{proposition}[Closure of executable support]
\label{prop:closure}
For any $x\in\X$ and $a\in\A(x)$, the executor guards imply
\begin{equation*}
\T(x,a)\in\X.
\end{equation*}
\end{proposition}

Consequently, every state reached by a stochastic process supported on $\{\A(x)\}_{x\in\X}$ remains in the declared molecular state space. The proof is given in \cref{app:proofs}.

The legal edit sets $\A(x)$ therefore determine which molecular transitions are possible, but not how probability should be distributed among them. We next learn a goal-independent stochastic law over this executable support.

\subsection{Learned molecular reference process}
\label{sec:rtheta}

Having defined which molecular rewrites are executable, we next learn how probability is distributed over them. For every state $x$ with $\A(x)\neq\emptyset$, the reference model defines a normalized conditional law over legal edit marks,
\begin{equation}
p_\theta(a\mid x,t),\qquad\sum_{a\in\A(x)}p_\theta(a\mid x,t)=1,
\label{eq:marklaw}
\end{equation}
where $t$ is a progress variable. Writing an edit as $a=(k,o,\eta)$, with family $k$, matched molecular operands $o$, and typed payload $\eta$, we parameterize
\begin{equation}
p_\theta(a\mid x,t)=p_\theta(k\mid x,t)\,p_\theta(o,\eta\mid k,x,t).
\label{eq:hierarchical}
\end{equation}
A shared molecular encoder represents the current graph, while family-specific scorers normalize only over candidates admitted by the executor. The reference model receives no downstream objective, remaining budget, or task-specific constraint, and the original source molecule is not supplied as an additional conditioning variable. The source determines the initial state; task information enters later through support restrictions and control. Thus the learned reference law represents molecular \emph{plausibility} independently of molecular \emph{purpose}.

Supervision is obtained by compiling molecular endpoint pairs into executable edit programs. Every stored teacher transition is replayed through the same executor used at inference, so the model is trained only on transitions available to the deployed process. The compiler may use both endpoints to construct an edit program, but neither the target endpoint nor the program is available to the reference model at inference. Because multiple legal marks can execute to the same molecular graph, the learning objective is defined on the resulting molecular successor rather than on a particular syntactic realization. Let $R_{\theta,t}(y\mid x)$ denote the induced law over distinct committed molecular successors, formalized in \cref{sec:canonical}. For productive teacher transitions $x_i\rightarrow y^{\star}_i$, we minimize
\begin{equation}
\mathcal{L}_{\mathrm{ref}}(\theta)=-\frac{\sum_i w_i\log R_{\theta,t_i}(y^{\star}_i\mid x_i)}{\sum_i w_i},
\label{eq:refloss}
\end{equation}
where $w_i$ is the registered weight under the declared editing measure. Terminal rows do not enter this objective. \Cref{eq:refloss} therefore learns which committed molecular successor follows a state under that measure; it does not calibrate continuous-time holding rates.

This objective is related to generator matching but is not the generator-matching population objective. Writing the productive rate as $Q^{+}(y\mid x,t)=\Lambda(x,t)P(y\mid x,t)$, the corresponding generator-matching divergence separates pointwise into a productive-hazard term and a conditional-successor term proportional to $\Lambda^{\star}\KL(P^{\star}\,\|\,R_\theta)$. \method{} optimizes the successor component only: the hazard term is not optimized, and \cref{eq:refloss} is averaged under the declared family-weighted editing measure rather than the state--time marginal of the conditional-process mixture \citep{holderrieth2025generator}. \Cref{app:gm} gives the full decomposition and the precise measure-level distinction.

The learned mark law is not yet the molecular kernel used for control: distinct executable marks can yield the same molecular state. We therefore next push $p_\theta$ through the executor and aggregate probability by canonical molecular outcome.

\subsection{Canonical molecular successors}
\label{sec:canonical}

We now push the learned edit law through the executor to obtain a stochastic process over distinct molecular states. Multiple legal marks can produce the same molecular graph because of graph symmetries, equivalent matches, or alternative internal rewrite descriptions. For a canonical successor $y$, define its successor fiber
\begin{equation}
\Gset_y(x)=\bigl\{\,a\in\A(x):\ \T(x,a)\simeq y\,\bigr\},
\label{eq:sucfiber}
\end{equation}
where $\simeq$ denotes equivalence under the registered molecular canonicalization. The induced probability of $y$ is
\begin{equation}
\bar p_{\theta,t}(y\mid x)=\sum_{a\in\Gset_y(x)}p_\theta(a\mid x,t).
\label{eq:pushforward}
\end{equation}
Thus probability assigned to a molecular transition is the aggregate mass of every executable mark that realizes it. For example, deleting any of the three symmetry-equivalent terminal methyl carbons in 2-methylpropane, $\op{CC(C)C}$, produces propane; canonicalization exposes propane once with the combined probability of those marks rather than overweighting it because of representational multiplicity.

Fixed-budget editing counts committed changes of molecular state rather than internal mark draws. Let $\Sset(x)$ denote the distinct non-self canonical successors admitted by this convention. Conditioning the pushforward distribution on a committed state change defines the molecular reference kernel
\begin{equation}
R_{\theta,t}(y\mid x)=\frac{\one\{y\in\Sset(x)\}\,\bar p_{\theta,t}(y\mid x)}{\sum_{y'\in\Sset(x)}\bar p_{\theta,t}(y'\mid x)}.
\label{eq:rtheta}
\end{equation}
If $\Sset(x)=\emptyset$, $x$ is terminal for the committed molecular process. For $b>0$, such a state has $h_b(x,z)=0$ under the backward recursion of \cref{sec:control}; trajectories terminating before the declared horizon therefore carry zero mass in the exact-$K$-step controlled terminal law. Otherwise $R_{\theta,t}$ is a normalized law over distinct molecular successors, and iterating it defines the goal-independent reference process optimized by \cref{eq:refloss}. We suppress $t$ and write $R_\theta(y\mid x)$ when the registered progress convention is fixed; canonical self-events and other non-committed outcomes are detailed in \cref{app:quotient}.

\begin{proposition}[Canonical normalization and refinement invariance]
\label{prop:canonical}
For every nonterminal state $x$, $\sum_{y\in\Sset(x)}R_{\theta,t}(y\mid x)=1$. Moreover, any refinement or coarsening of the internal mark representation that preserves the aggregate probability mass of every canonical successor fiber leaves $R_{\theta,t}$ unchanged. Consequently, state-level controllers that reweight $R_{\theta,t}$ using quantities defined on canonical molecular successors are invariant to such representational refinements.
\end{proposition}

This invariance relies on performing control after canonicalization: reweighting or truncating individual marks beforehand can reintroduce dependence on arbitrary edit multiplicities. With the goal-independent reference process now defined on complete molecular states, we next control this frozen kernel toward finite-horizon design objectives.
